\documentclass[11pt,a4paper]{article}

% =========================================================
% Essential packages (same style as your example)
% =========================================================
\usepackage[utf8]{inputenc}
\usepackage[T1]{fontenc}
\usepackage[english]{babel}
\usepackage{amsmath,amssymb,amsfonts}
\usepackage{geometry}
\usepackage{xcolor}
\usepackage{titlesec}
\usepackage{tcolorbox}
\usepackage{fancyhdr}
\usepackage{enumitem}
\usepackage{booktabs}
\usepackage{array}
\usepackage{microtype}
\usepackage{hyperref}
\usepackage{tabularx}

% =========================================================
% Page geometry
% =========================================================
\geometry{
top=2.5cm,
bottom=2.5cm,
left=2.5cm,
right=2.5cm,
headheight=14pt
}

% =========================================================
% Color definitions
% =========================================================
\definecolor{darkblue}{RGB}{30,58,138}
\definecolor{mediumblue}{RGB}{59,130,246}
\definecolor{lightblue}{RGB}{239,246,255}
\definecolor{grayblue}{RGB}{241,245,249}

% =========================================================
% Header and footer
% =========================================================
\pagestyle{fancy}
\fancyhf{}
\fancyhead[L]{\textcolor{darkblue}{\textbf{Trusted Scalar}}}
\fancyhead[R]{\textcolor{darkblue}{\thepage}}
\renewcommand{\headrulewidth}{1pt}
\renewcommand{\headrule}{\hbox to\headwidth{\color{darkblue}\leaders\hrule height \headrulewidth\hfill}}

% =========================================================
% Section formatting
% =========================================================
\titleformat{\section}
{\Large\bfseries\color{darkblue}}
{\thesection}{1em}{}
[\vspace{-0.5em}\textcolor{darkblue}{\hrule height 2pt}]

\titleformat{\subsection}
{\large\bfseries\color{darkblue}}
{\thesubsection}{1em}{}
[\vspace{-0.5em}\textcolor{mediumblue}{\hrule height 1pt}]

\titleformat{\subsubsection}
{\normalsize\bfseries\color{darkblue}}
{\thesubsubsection}{1em}{}

% =========================================================
% Custom environments
% =========================================================
\newtcolorbox{definition}{
colback=lightblue,
colframe=darkblue,
fonttitle=\bfseries,
left=5pt,
right=5pt,
top=5pt,
bottom=5pt,
arc=3pt,
boxrule=0.5pt
}
\newtcolorbox{schema}{
colback=grayblue,
colframe=darkblue,
fonttitle=\bfseries,
left=10pt,
right=10pt,
top=8pt,
bottom=8pt,
arc=0pt,
boxrule=0pt,
leftrule=4pt
}

% =========================================================
% Hyperref setup
% =========================================================
\hypersetup{
colorlinks=true,
linkcolor=darkblue,
urlcolor=darkblue,
citecolor=darkblue
}

% =========================================================
% List formatting
% =========================================================
\setlist[itemize]{leftmargin=1cm}
\setlist[enumerate]{leftmargin=1cm}

% =========================================================
% Document begins
% =========================================================
\begin{document}

% =========================================================
% Title page
% =========================================================
\begin{titlepage}
\centering
\vspace*{2cm}

{\Huge\bfseries\color{darkblue} Trusted Scalar}\\[0.5cm]
{\Large\color{darkblue} ROM Qualification as a Safety-Critical Scalar Contract}\\[2cm]

{\large\itshape A rigorous formulation for treating reduced-order model outputs\\
as certificate-like quantities inside safety and invariance control laws}\\[4cm]

\vfill
{\large\color{darkblue}\today}

\end{titlepage}

% Table of contents
\tableofcontents
\newpage

% =========================================================
% Passage (verbatim content, not edited)
% =========================================================

\begin{definition}
\textbf{Trusted Scalar (core claim)}\\
\textbf{a reduced-order model (ROM) with validated predictors can serve as a “trusted scalar”} \textit{if} you qualify it the same way you’d qualify a \textbf{certificate, margin, or invariant index}: with \textbf{bounded error}, \textbf{integrity}, and \textbf{contractual meaning} inside the control law.

In your thesis, this is already the role played by the composite index \textbf{$\varepsilon$} and the robust margin \textbf{$\delta(\varepsilon)$} in the invariant + CBF recursion.

Below is the rigorous way to qualify ROM outputs as “trusted scalars.”
\end{definition}

\hrule
\vspace{0.75em}

\section{What a ``trusted scalar'' actually means (rigorously)}

A scalar (s\_t) is \textbf{trusted} if it satisfies a \textbf{contract} of the form:

\subsection*{\textbf{(C1) Semantic contract}}

It must \textit{mean} something safety-relevant, not just correlate.

Example contracts:

\begin{itemize}
\item (s\_t = $\varepsilon_t$) is a \textbf{risk index} driving conservatism
\item (s\_t = $\delta(\varepsilon_t)$) is a \textbf{robust safety margin} ensuring forward invariance
\end{itemize}

\subsection*{\textbf{(C2) Boundedness contract}}

It must be \textbf{bounded and well-formed}:
\[
s_t \in [s_{\min}, s_{\max}], \quad \text{no NaNs, no undefined states}
\]

\subsection*{\textbf{(C3) Conservatism contract}}

If it’s wrong, it must fail \textbf{in the conservative direction}.

That is the key rule:

\begin{schema}
\textbf{Conservatism rule}\\[-0.25em]
\begin{quote}
\textbf{Overestimate risk is acceptable; underestimate risk is not.}
\end{quote}
\end{schema}

\hrule
\vspace{0.75em}

\section{How ROM predictors become ``trusted scalars''}

A ROM doesn’t become trusted because it’s ``accurate on average.''

It becomes trusted because you can \textbf{upper-bound its error} in the operational envelope and \textbf{compile that bound into the safety margin}.

Your own framework explicitly assumes bounded prediction/encoding errors:

\begin{itemize}
\item (f,g) Lipschitz
\item ($|\Delta f|\le \epsilon_f$), ($|\Delta g|\le \epsilon_g$)
\item encoder error bounded by ($\epsilon_z$)
\end{itemize}

And then it gives a \textbf{sufficient condition} for safety:

If
\[
\delta(\varepsilon)\ge \text{(terms that dominate worst-case deficits from } \epsilon_f,\epsilon_g,\epsilon_z\text{)}
\]
then the CBF-filtered action guarantees forward invariance .

That is literally a \textit{qualification recipe}:

\begin{schema}
\textbf{Qualification recipe}\\[-0.25em]
\begin{quote}
ROM predictors are trusted \textbf{only through the margin} that dominates their worst-case error.
\end{quote}
\end{schema}

So the ROM doesn’t need to be perfect.
It needs to be \textbf{bounded} and \textbf{accounted for}.

\hrule
\vspace{0.75em}

\section{Qualification pipeline: ROM $\to$ trusted scalar (step-by-step)}

\subsection*{Step A — Define the scalar’s safety contract}

Pick which scalar the ROM will output:

\subsubsection*{Option 1: ROM outputs \textbf{$\varepsilon$} (risk index)}

You already define:
\[
\varepsilon = \text{median}(|\xi|)+\alpha,\text{Var}(sct)+\beta,\mathbf 1[S=1]
\]

This scalar is \textit{allowed} to be approximate \textbf{as long as error is conservative} ($\varepsilon$ too high $\Rightarrow$ more conservatism).

\subsubsection*{Option 2: ROM outputs \textbf{$\delta$} (robust margin)}

This is stronger: $\delta$ directly participates in the safety inequality .

If $\delta$ is the trusted scalar, it must satisfy:
\[
\delta_{\text{ROM}}(t)\;\ge\;\delta_{\text{required}}(t)
\]
for all operating points in the certified envelope.

\textbf{This is a certificate-style requirement.}

\hrule
\vspace{0.75em}

\subsection*{Step B — Bound ROM prediction error \textit{in the relevant norm}}

You qualify the predictor by establishing a bound:
\[
|y_{\text{true}} - y_{\text{ROM}}| \le b(x)
\]

Then you simplify it into a worst-case constant bound in the envelope:
\[
|y_{\text{true}} - y_{\text{ROM}}| \le \bar b
\]

This is how your ($\epsilon_f,\epsilon_g,\epsilon_z$) enter the proof chain.

\textbf{Key:} you’re not proving ``high accuracy.''
You’re proving \textbf{upper-bounded deviation}.

\hrule
\vspace{0.75em}

\subsection*{Step C — Convert predictor error into a conservative scalar mapping}

You then define the trusted scalar as a monotone function of the bound:

\subsubsection*{If scalar is $\delta$:}

\[
\delta(\varepsilon) \triangleq \delta_0 + k\cdot \varepsilon
\]
with (k) chosen so $\delta$ dominates the error terms.

Your Proposition is exactly the form: $\delta$ must exceed a sum of worst-case deficits .

\subsubsection*{If scalar is $\varepsilon$:}

you ensure \textbf{$\varepsilon$ underestimation is bounded}, and add a safety bias:
\[
\varepsilon_{\text{trusted}} = \varepsilon_{\text{ROM}} + \Delta_{\text{bias}}
\]
with ($\Delta_{\text{bias}}$) chosen so:
\[
\varepsilon_{\text{trusted}} \ge \varepsilon_{\text{true}}
\]
over the certified envelope.

\hrule
\vspace{0.75em}

\subsection*{Step D — Add \textit{integrity} and \textit{anti-spoofing} qualification}

A scalar can be mathematically valid but not \textbf{trustworthy} if software can tamper with it.

So you qualify the signal path:

\begin{itemize}
\item computed in an isolated block
\item bounds checked + saturating
\item signed/hashed or bus-authenticated
\item repeated by redundancy if high stakes
\end{itemize}

If the scalar is corrupted, it must default to \textbf{max risk}.

That creates \textit{fail-closed trust}.

\hrule
\vspace{0.75em}

\subsection*{Step E — Runtime qualification (online conformance monitoring)}

Even a qualified ROM can become invalid outside its envelope.

So you enforce ``trust only while conforming.''

Examples:

\begin{itemize}
\item residual tests: ($|y - \hat y|$)
\item drift / non-equilibrium detection (your (sct) explicitly captures this)
\item mode flag (S = 1[n\_z $\ne$ n\_p]) is a structural anomaly indicator
\end{itemize}

If conformance fails $\to$ scalar escalates conservatism:

\begin{itemize}
\item $\varepsilon$ increases
\item $\delta$ inflates
\item trust region shrinks (you already specify this behavior)
\end{itemize}

This is an extremely clean principle:

\begin{schema}
\textbf{Runtime qualification principle}\\[-0.25em]
\begin{quote}
\textbf{Qualification is not a one-time stamp. It is continuously re-earned.}
\end{quote}
\end{schema}

\hrule
\vspace{0.75em}

\section{The rigorous reason this works: it composes with your forward-invariance proof}

Your stack has the exact structure needed:

\begin{enumerate}
\item \textbf{IS--MPC constraints} with ($X_{\text{inv}}(\varepsilon)$) shrinking with $\varepsilon$
\item A \textbf{CBF-QP safety filter} that projects nominal actions into safe ones
\item An explicit \textbf{robust margin} $\delta(\varepsilon)$ sized to dominate modeling/prediction error
\item Online schedules that respond to $\varepsilon$ by shrinking trust radius and inflating $\delta$ when infeasible
\end{enumerate}

So a ROM scalar becomes ``trusted'' precisely because:

\begin{itemize}
\item it \textbf{does not need to be correct}
\item it needs to be \textbf{conservative and bounded}
\item and its conservatism is what \textbf{preserves invariance}.
\end{itemize}

That’s a real safety argument, not vibes.

\hrule
\vspace{0.75em}

\section{Practical qualification criteria (what you can put in a spec)}

A ROM-derived scalar (s\_t) is \textit{qualified as trusted} if:

\subsection*{\textbf{Q1 — Bounded error}}

There exists a certified envelope (D) such that:
\[
\sup_{x\in D}|y_{\text{true}}-y_{\text{ROM}}|\le \bar b
\]

\subsection*{\textbf{Q2 — Conservative mapping}}

The scalar is biased so that it never reduces conservatism when wrong.

\subsection*{\textbf{Q3 — Contract integration}}

The scalar appears only inside monotone ``safety tightening'' mechanisms:

\begin{itemize}
\item shrinking feasible sets
\item inflating $\delta$ margins
\item tightening actuation clamps
  (not loosening them)
\end{itemize}

\subsection*{\textbf{Q4 — Online conformance}}

When conformance tests fail, the scalar moves to \textbf{max risk} behavior.

\subsection*{\textbf{Q5 — Integrity}}

Scalar path is isolated, bounds-checked, and non-bypassable.

\hrule
\vspace{0.75em}

\section*{Bottom line}

\textbf{Yes} — ROM predictors can be ``trusted scalars,'' but \textit{only} after qualification proves:

\begin{itemize}
\item bounded error in the operating envelope
\item conservative failure direction
\item integration through $\delta/\varepsilon$ monotone safety tightening
\item continuous runtime conformance
\item enforcement integrity
\end{itemize}

\end{document}
